\documentclass[11pt, letterpaper]{article}

% ---- Packages ----
\usepackage[margin=0.75in]{geometry}
\usepackage{titlesec}
\usepackage{enumitem}
\usepackage{hyperref}
\usepackage{xcolor}
\usepackage{fontenc}
\usepackage[utf8]{inputenc}
\usepackage{mathpazo} % Palatino font
\usepackage{microtype}
\usepackage{fancyhdr}
\usepackage{lastpage}
\usepackage{tcolorbox}
\usepackage{tabularx}

% ---- Colors ----
\definecolor{accent}{HTML}{1A3C6E}
\definecolor{linkblue}{HTML}{2B5797}
\definecolor{highlight}{HTML}{FFF8E1}
\definecolor{highlightborder}{HTML}{E8A817}
\definecolor{gray}{HTML}{666666}

% ---- Hyperlinks ----
\hypersetup{
  colorlinks=true,
  linkcolor=linkblue,
  urlcolor=linkblue,
  pdfauthor={Minxing Pang},
  pdftitle={Minxing Pang - Curriculum Vitae},
}

% ---- Section Formatting ----
\titleformat{\section}
  {\large\bfseries\color{accent}\scshape}
  {}
  {0em}
  {}
  [\vspace{-0.6em}\textcolor{accent}{\rule{\textwidth}{0.8pt}}]
\titlespacing*{\section}{0pt}{14pt}{6pt}

% ---- Footer ----
\pagestyle{fancy}
\fancyhf{}
\renewcommand{\headrulewidth}{0pt}
\fancyfoot[C]{\small\textcolor{gray}{Minxing Pang --- Curriculum Vitae --- Page \thepage\ of \pageref{LastPage}}}

% ---- Custom commands ----
\newcommand{\cventry}[2]{%
  \noindent\begin{tabularx}{\textwidth}{@{}X r@{}}
    #1 & \textcolor{gray}{#2} \\
  \end{tabularx}\vspace{1pt}%
}

\newcommand{\me}[1]{\textbf{#1}}

\newcommand{\papertag}[2]{\,[\href{#2}{\textcolor{linkblue}{#1}}]}

% ---- Highlight box for featured paper ----
\tcbset{
  highlightbox/.style={
    colback=highlight,
    colframe=highlightborder,
    boxrule=0.6pt,
    arc=2pt,
    left=6pt,
    right=6pt,
    top=4pt,
    bottom=4pt,
    boxsep=0pt,
  }
}

% ================= DOCUMENT =================
\begin{document}

\thispagestyle{fancy}

% ---- Header ----
\begin{center}
  {\LARGE\bfseries\color{accent}\textls[120]{MINXING PANG}}\\[6pt]
  {\small
    \href{mailto:minxing@sas.upenn.edu}{minxing@sas.upenn.edu} \\[2pt]
    \href{https://github.com/maxpmx}{GitHub} \enspace$\vert$\enspace
    \href{https://scholar.google.com/citations?user=wZbrjsQAAAAJ}{Google Scholar}
  }
\end{center}

\vspace{-4pt}

% =========================================================
\section{Research Interests}
% =========================================================

AI for Healthcare\enspace$\cdot$\enspace Biomedical Imaging\enspace$\cdot$\enspace Spatial Omics\enspace$\cdot$\enspace Computational Pathology\enspace$\cdot$\enspace Deep Learning

\vspace{4pt}
\noindent\small\textit{My research develops deep learning methods that bridge routine histology imaging and molecular measurement, turning ubiquitous tissue images into rich, spatially resolved, multi-modal readouts for precision medicine.}

% =========================================================
\section{Education}
% =========================================================

\cventry{\textbf{University of Pennsylvania}, Philadelphia, PA, USA}{2022 -- 2026 (Expected)}
\hspace{1em}\textit{Ph.D.\ in Applied Mathematics and Computational Science}\\
\hspace{1em}Advisor: Kai Tan, Ph.D., Professor of Pediatrics

\vspace{6pt}
\cventry{\textbf{King Abdullah University of Science and Technology (KAUST)}, Saudi Arabia}{2020 -- 2022}
\hspace{1em}\textit{M.S.\ in Applied Mathematics and Computational Science}\\
\hspace{1em}Advisor: Jesper Tegn\'er, Ph.D., Professor of Bioscience

\vspace{6pt}
\cventry{\textbf{Nankai University}, Tianjin, China}{2016 -- 2020}
\hspace{1em}\textit{B.S.\ in Information and Computational Science}\\
\hspace{1em}Advisor: Jishou Ruan, Ph.D., Professor of Mathematics\\
\hspace{1em}Thesis: \textit{Modeling single-cell RNA-seq data using deep learning}

% =========================================================
\section{Publications}
% =========================================================

\noindent\small\textcolor{gray}{Name in \textbf{bold} indicates author. Google Scholar profile linked in header.}\normalsize
\vspace{4pt}

\begin{enumerate}[leftmargin=1.5em, itemsep=6pt, label={[\arabic*]}]

\item \me{Pang M}, Roy TK, Wu X, Tan K. (2025). CelloType: A Unified Model for Segmentation and Classification of Tissue Images. \textit{Nature Methods}.
\papertag{Paper}{https://www.nature.com/articles/s41592-024-02513-1}
\papertag{Codes}{https://github.com/tanlabcode/CelloType}

\item Bandyopadhyay S, Duffy M, Ahn K, Sussman J, \me{Pang M}, \ldots, and Tan K. (2024). Mapping the cellular biogeography of human bone marrow niches using single-cell transcriptomics and proteomic imaging. \textit{Cell}.
\papertag{Paper}{https://www.cell.com/cell/fulltext/S0092-8674(24)00408-2}

\item Yu W, Biyik-Sit R, Uzun Y, Chen CH, Thadi A, Sussman JH, \me{Pang M}, \ldots, and Tan K. (2025). Longitudinal single-cell multiomic atlas of high-risk neuroblastoma reveals chemotherapy-induced tumor microenvironment rewiring. \textit{Nature Genetics}.

\end{enumerate}

% =========================================================
\section{Conference Papers}
% =========================================================

\begin{enumerate}[leftmargin=1.5em, itemsep=6pt, label={[\arabic*]}, resume]

\item \me{Pang M}, Tegn\'er J. (2020). Representation Learning and Translation between the Mouse and Human Brain using a Deep Transformer Architecture. \textit{International Conference on Machine Learning (ICML) 2020 Workshop on Computational Biology}. \textbf{Contributed Talk} (Top 5\%).
\papertag{PDF}{https://icml-compbio.github.io/icml-website-2020/2020/papers/WCBICML2020_paper_29.pdf}
\papertag{Codes}{https://github.com/maxpmx/scTT-pytorch}

\end{enumerate}

% =========================================================
\section{Manuscripts Under Review / Preprints}
% =========================================================

\begin{enumerate}[leftmargin=1.5em, itemsep=6pt, label={[\arabic*]}, resume]

\item
% \begin{tcolorbox}[highlightbox]
\me{Pang M}, Roy TK, Wu X, Tan K. (2025). A Query-driven generative AI synthesizes multi-modal spatial omics from histology. \textbf{\textit{Under review}}.
\papertag{bioRxiv}{https://www.biorxiv.org/content/10.64898/2025.12.11.693669v1}
% \hfill\textcolor{highlightborder}{\small$\bigstar$ \textbf{Featured Work}}
% \end{tcolorbox}

\item \me{Pang M}, Su K, Li M. (2021). Leveraging information in spatial transcriptomics to predict super-resolution gene expression from histology images in tumors.
\papertag{bioRxiv}{https://www.biorxiv.org/content/10.1101/2021.02.28.433265}
\papertag{Codes}{https://github.com/maxpmx/HisToGene}

% \item \me{Pang M}, Tegn\'er J. (2020). Multitask learning for Transformers with application to large-scale single-cell transcriptomes.
% \papertag{bioRxiv}{https://www.biorxiv.org/content/10.1101/2021.05.29.445901}

\end{enumerate}

% =========================================================
\section{Invited Talks \& Presentations}
% =========================================================

\textit{Towards Unified Segmentation and Classification of Tissue Images}

\vspace{4pt}
\cventry{\hspace{1em}HTAN-HuBMAP Joint Meeting, Stanford University, CA}{March 2024}
\cventry{\hspace{1em}Schmidt Center Symposium: Biomedical Science and AI, Koch Institute, MIT, MA}{April 2025}

\vspace{6pt}
\noindent \textit{Generative AI synthesizes multi-modal spatial omics from histology}

\vspace{4pt}
\cventry{\hspace{1em}Shanghai Medical College, Fudan University, CN \quad\textcolor{gray}{(Invited)}}{December 2024}
\cventry{\hspace{1em}School of Medicine, Zhejiang University, CN \quad\textcolor{gray}{(Invited)}}{July 2025}


% =========================================================
\section{Honors \& Awards}
% =========================================================

% \cventry{National Scholarship for Outstanding International Students, China Scholarship Council}{2025}
\cventry{KAUST Fellowship}{2020 -- 2022}
\cventry{Outstanding Undergraduate Thesis of Nankai University (Top 3\%)}{2020}
\cventry{Nankai University Innovation Scholarship}{2017, 2018}
\cventry{Provincial First Prize in Chinese Physics Competition}{2015}

% =========================================================
\section{Mentoring \& Teaching}
% =========================================================

\cventry{\textbf{Research Mentor}, Tan Lab, Children's Hospital of Philadelphia \& UPenn}{2022 -- Present}
\hspace{1em}Mentor junior graduate and undergraduate students on projects in computational pathology and spatial omics,\par
\hspace{1em}guiding method development, coding, and manuscript preparation.

% =========================================================
\section{Technical Skills}
% =========================================================

\textbf{Advanced:} Python, PyTorch, Deep Learning, Biomedical Imaging Analysis\\
\textbf{Familiar:} R, C++

% =========================================================
\section{References}
% =========================================================

References available upon request.

\end{document}
